\section{Ethical, Societal and Sustainability Considerations}
\textbf{Privacy by design.} No real patient data is used; every encounter is synthetic and labelled so.
Direct identifiers stop at bronze: silver and gold hold only the keyed hash and age bands, the key never enters
Git, and every export hides counts from 1 to 4 with complementary suppression, so a total cannot reveal a hidden
cell. One residual risk is known and written in the code: the padded HMAC key appears in Spark's query plan, so
Spark event logs of the silver job must be treated as secret and kept outside the repository.

\textbf{Record-level extracts and public release (PRV-4).} \texttt{pipeline/privacy/kanon.py} measures
k-anonymity~\cite{sweeney2002k} over the quasi-identifiers (age band, sex, district, onset week, disease,
facility) of any record-level extract. k-anonymity does not stop a homogeneity attack, where everyone in a class
shares one diagnosis, so a disease extract also needs l-diversity. For a public release, differential
privacy~\cite{dwork2006dp} is the stronger tool: the Laplace mechanism adds noise with standard deviation
$\sqrt{2}/\varepsilon$ to each count, about 1.4 cases at $\varepsilon=1$. The privacy budget adds up over every
published cell and week, and noisy counts can be negative and must be post-processed. We would use it for the
public release only; ministry users keep exact counts with suppression.
\begin{teammate}{Samira: confirm}
Equity (what $\alpha$ protects and what it costs), misuse (stigma of a named district, alert fatigue), the
limits of big-data signals~\cite{lazer2014flu}, and the energy and hardware footprint.
\end{teammate}
